\documentclass[12pt, a4paper]{article}
\usepackage[utf8]{inputenc}
\usepackage[T1]{fontenc}
\usepackage{amsmath, amssymb, amsthm}
\usepackage{geometry}
\usepackage{listings}
\geometry{margin=2.5cm}

\newtheorem{theorem}{Theorem}[section]
\newtheorem{lemma}[theorem]{Lemma}
\newtheorem{definition}[theorem]{Definition}

\title{On the Erdős-Straus Conjecture: Algebraic Analysis and Modular Decomposition}
\author{Charles EDOU NZE\thanks{Charles EDOU NZE, chercheur indépendant}}
\date{\today}

\lstdefinelanguage{lean}{
  keywords={import, def, theorem, lemma, by, sorry, Prop, Nat, open, section, Exists, fun},
  sensitive=true,
  comment=[l]--
}

\begin{document}
\maketitle

\begin{abstract}
This document presents a rigorous analysis and structural decomposition of the Erdős-Straus conjecture. We lay out axiomatic definitions, review relevant recent literature, isolate key lemmas for specific congruence classes, and provide an architecture for autoformalization in a proof assistant like Lean 4.
\end{abstract}

\tableofcontents
\newpage

\section{Axiomatic Definitions and Context}
\begin{definition}
For every integer $n \in \mathbb{N}$ with $n \ge 2$, the Erdős-Straus equation is defined as the Diophantine equation:
\[\frac{4}{n} = \frac{1}{x} + \frac{1}{y} + \frac{1}{z}\]
where $x, y, z \in \mathbb{Z}_{>0}$.
\end{definition}

\subsection{Contextual Literature Research}
The problem, formulated by Paul Erdős and Ernst G. Straus in 1948, posits that the equation always has a positive integer solution for $n \ge 2$. Recent studies by authors such as Miguel Angel Lopez have classified solutions into types, such as Type A and Type B, by defining a complete congruence system. Furthermore, Philemon Urbain Mballa has explored an unexpected connection between the discrete zeta function and the Erdős-Straus conjecture through additive decomposition. This methodology shows strong analogies to the algebraic manipulations used in solving Pell-Fermat equations, where solutions are grouped by modular invariants and analyzed via infinite descents or structural classifications.

\section{Proof Strategy and Lemma Isolation}
The fundamental approach involves partitioning the integers $n$ by their congruence class modulo an integer, historically $840$, although general cases can be studied for primes.

\begin{lemma}
For $n = 4k+3$ with $k \in \mathbb{N}$, the equation admits a positive integer solution.
\end{lemma}
\begin{proof}
Let $n = 4k+3$.
We set $x = (n+1)/4 = (4k+4)/4 = k+1$. Because $n \ge 3$, $k \ge 0$, and thus $x \ge 1$, ensuring $x \in \mathbb{Z}_{>0}$.
We evaluate the remaining difference:
\begin{align*}
\frac{4}{n} - \frac{1}{x} &= \frac{4}{n} - \frac{1}{k+1} \\
&= \frac{4(k+1) - n}{n(k+1)} \\
&= \frac{4k+4 - (4k+3)}{n(k+1)} \\
&= \frac{1}{n(k+1)}
\end{align*}
The remaining fraction $\frac{1}{n(k+1)}$ must be split into two unit fractions $\frac{1}{y} + \frac{1}{z}$.
We employ the standard Egyptian fraction splitting identity:
\[ \frac{1}{A} = \frac{1}{A+1} + \frac{1}{A(A+1)} \]
Applying this to $A = n(k+1)$, we set:
$y = n(k+1) + 1$
and
$z = n(k+1)(n(k+1)+1)$
Since $n \ge 3$ and $k \ge 0$, both $y$ and $z$ are strictly positive integers.
By substitution:
\begin{align*}
\frac{1}{x} + \frac{1}{y} + \frac{1}{z} &= \frac{1}{k+1} + \frac{1}{n(k+1) + 1} + \frac{1}{n(k+1)(n(k+1)+1)} \\
&= \frac{1}{k+1} + \frac{n(k+1)+1 + 1}{(n(k+1)+1)(n(k+1)(n(k+1)+1))} \\
\end{align*}
Wait, the direct sum of the last two is exactly $\frac{1}{n(k+1)}$. Let us show it explicitly:
\begin{align*}
\frac{1}{y} + \frac{1}{z} &= \frac{1}{n(k+1)+1} + \frac{1}{n(k+1)(n(k+1)+1)} \\
&= \frac{n(k+1)}{n(k+1)(n(k+1)+1)} + \frac{1}{n(k+1)(n(k+1)+1)} \\
&= \frac{n(k+1)+1}{n(k+1)(n(k+1)+1)} \\
&= \frac{1}{n(k+1)}
\end{align*}
Adding $\frac{1}{x} = \frac{1}{k+1}$, we get:
\begin{align*}
\frac{1}{x} + \frac{1}{y} + \frac{1}{z} &= \frac{1}{k+1} + \frac{1}{n(k+1)} \\
&= \frac{n}{n(k+1)} + \frac{1}{n(k+1)} \\
&= \frac{n+1}{n(k+1)} \\
&= \frac{4k+4}{n(k+1)} \\
&= \frac{4(k+1)}{n(k+1)} \\
&= \frac{4}{n}
\end{align*}
This proves the lemma for all $n \equiv 3 \pmod 4$.
\end{proof}

\section{Architecture for Autoformalization}
The formalization of the Erdős-Straus conjecture requires structuring the statement and the search space decomposition in Lean 4.

\begin{lstlisting}[language=lean, basicstyle=\ttfamily\small, breaklines=true]
import Mathlib.Data.Nat.Basic
import Mathlib.Tactic.Omega
import Mathlib.Tactic.Ring

-- Axiomatic definition of the Erdos-Straus property
def SatisfiesErdosStraus (n : Nat) : Prop :=
  Exists (fun x => Exists (fun y => Exists (fun z => x > 0 /\ y > 0 /\ z > 0 /\ 4 * x * y * z = n * (y * z + x * z + x * y))))

-- Lemma for mod 4 = 3
lemma erdos_straus_mod_4_3 (k : Nat) : SatisfiesErdosStraus (4 * k + 3) := by
  sorry

-- Main Theorem
theorem erdos_straus_conjecture (n : Nat) (hn : n >= 2) : SatisfiesErdosStraus n := by
  sorry
\end{lstlisting}

\end{document}
